\documentclass[12pt, a4paper, french]{article}

% ========================
% IMPORTATION DES PACKAGES
% ========================

\usepackage[T1]{fontenc}
\usepackage[utf8]{inputenc}
\usepackage[margin=2cm]{geometry}
\usepackage{lmodern}
\usepackage{theorem}
\usepackage{babel}

\pagestyle{headings}

\renewcommand{\familydefault}{\sfdefault}

\theoremstyle{break}
\theoremheaderfont{\scshape}
\theorembodyfont{\upshape}

\newtheorem{exo}{Exercice}[section]
\newtheorem{sol}{Solution}[section]

\title{EXERCICE DE DIFFERENTIABILITÉ DANS UN BANACH}
\date{\today}
\author{}

\begin{document}
\section*{PRECISION}
Je tient à préciser que cette document est un recueille des notes d'exercice et de solution de cours mathématique intitulé la differentiabilite dans un Banach. Si cette support peut vous aider alors je vous invite à l'explorer. Tous les solutions et énoncé sont rediger selon mon propre formulation.

\section{Espace normé et espace de Banach}
\begin{exo}[Complétude]
       Montrer que les espaces suivant sont des espaces de Banach:
       \begin{itemize}
              \item[a)] $(\mathbb{R}^{n}, \| \cdot \|_{2})$
       \end{itemize}
\end{exo}

\begin{sol}
       \begin{itemize}
              \item[a)] Montrons que $(\mathbb{R}^{n}, \| \cdot \|_{2}) est un espace de Banach.$\\
              
       \end{itemize}
\end{sol}
\section{Application differentiable}


\section{Théorème des accroissement finis}


\section{Derivé d'ordre supérieur}


\section{Théorème d'inversion et implicite}
\end{document}
